% ID: SEL-D0013-BB-E03
% ESTADO: validada. S014; tres pasadas del mismo autor, sin segundo revisor.
\textbf{a) Dominio y ajuste en la unión.}\par\nopagebreak[4]
El dominio es $\mathbb R$. Aunque la primera expresión contiene $x-1$ en el denominador, se usa sólo para $x<0$; $x=1$ pertenece a la rama exponencial. Las dos ramas son continuas y derivables en sus interiores.
Para continuidad en $0$:
\[
2-a=a+b\quad\Longleftrightarrow\quad2=2a+b.
\]
Las derivadas de las ramas son $-a/(x-1)^2$ y $be^x$. Para igualar sus valores laterales en $0$, $-a=b$. Sustituyendo en la condición de continuidad:
\[
2=2a-a\quad\Longrightarrow\quad\boxed{a=2,\quad b=-2.}
\]
Con estos valores coinciden límites, valor de la función y derivadas laterales: $f(0)=0$ y $f'(0)=-2$.

\textbf{b) Monotonía y extremos relativos.}\par\nopagebreak[4]
\[
f'(x)=\begin{cases}-2/(x-1)^2&x<0,\\-2e^x&x>0,\end{cases}
\qquad f'(0)=-2.
\]
La derivada es estrictamente negativa en todo $\mathbb R$: \textbf{$f$ es estrictamente decreciente en todo su dominio y no tiene extremos relativos}. La unión no introduce un extremo.

\textbf{c) Asíntotas.}\par\nopagebreak[4]
\[
\lim_{x\to-\infty}\left(2+\frac2{x-1}\right)=2.
\]
Hay asíntota horizontal \boxed{y=2\text{ cuando }x\to-\infty}. Para $x\to+\infty$, $f(x)=2-2e^x\to-\infty$, sin asíntota horizontal. Tampoco hay oblicua: $f(x)/x\to-\infty$ a la derecha; a la izquierda la pendiente límite es cero y ya se obtiene la horizontal.
No hay asíntotas verticales, pues la función es continua y finita en todo punto real. En particular, $x=1$ no es una singularidad de la función por tramos.
